\chapter{Consolidated Paradigms}\label{app:paradigms}

The nominal tables give the established declensional classes and irregular
paradigms, followed by the verbal quantitative classes, aspect--alignment
exponents, reduced non-indicative opposition, and inner evidential forms.
Each table is reproduced from its principal treatment, where the conditions
on stem selection, case, alignment, and interpretation are specified.

\paradigmappendixtrue
\section{Nominal paradigms}

The principal treatments are given in Tables~\ref{tab:nm-plural-paradigm}, \ref{tab:city-town-paradigms}, \ref{tab:nm-paradigms-iv-vi}, \ref{tab:nm-paradigms-vii-ix}, and~\ref{tab:nm-irregular-paradigms}.

\begin{table}[ht]
  \footnotesize
  \caption{Complete paradigms of Classes I--III.}
  \ifparadigmappendix\else\label{tab:nm-plural-paradigm}\fi
  % Legacy labels for synoptic tables consolidated into the full paradigms.
  \ifparadigmappendix\else\label{tab:nm-stem-selection}\fi
  \ifparadigmappendix\else\label{tab:nm-def-allomorphy}\fi
  \ifparadigmappendix\else\label{tab:nm-def-pl-classiii}\fi
  \begin{tabularx}{\textwidth}{l l L L L L}
    \toprule
    {\sc noun (cl.)} & {\sc case} & {\sc sg} & {\sc pl}
      & {\sc sg.def} & {\sc pl.def} \\
    \midrule
    \ird{byl} (I) & \textsc{dir} & \ird{byl} & \ird{býle} & \ird{bylo} & \ird{býlete} \\
     & \textsc{trs} & \ird{byla} & \ird{býlov} & \ird{bylato} & \ird{býlove} \\
     & \textsc{ind} & \ird{bylam} & \ird{býlac} & \ird{bylamo} & \ird{býlace} \\
    \addlinespace
    \ird{lobek} (II) & \textsc{dir} & \ird{lobek} & \ird{lúbce} & \ird{lúbko} & \ird{lúbcé} \\
     & \textsc{trs} & \ird{lúbku} & \ird{lúbkov} & \ird{lúbkuto} & \ird{lúbkove} \\
     & \textsc{ind} & \ird{lúbcí} & \ird{lúbcem} & \ird{lúbcíte} & \ird{lúbceme} \\
    \addlinespace
    \ird{perek} (II) & \textsc{dir} & \ird{perek} & \ird{pérce} & \ird{pérko} & \ird{pércé} \\
     & \textsc{trs} & \ird{pérku} & \ird{pérkov} & \ird{pérkuto} & \ird{pérkove} \\
     & \textsc{ind} & \ird{pérce} & \ird{pércem} & \ird{pércé} & \ird{pérceme} \\
    \addlinespace
    \ird{trava} (III) & \textsc{dir} & \ird{trava} & \ird{trávy} & \ird{travó} & \ird{trávyte} \\
     & \textsc{trs} & \ird{travy} & \ird{tráv} & \ird{travyte} & \ird{tráve} \\
     & \textsc{ind} & \ird{travě} & \ird{trávam} & \ird{travěte} & \ird{trávame} \\
    \addlinespace
    \ird{stráz} (III) & \textsc{dir} & \ird{stráz} & \ird{strazy} & \ird{stráze} & \ird{strazyte} \\
     & \textsc{trs} & \ird{strazy} & \ird{stráz} & \ird{strazyte} & \ird{stráze} \\
     & \textsc{ind} & \ird{straže} & \ird{strazám} & \ird{stražete} & \ird{strazáme} \\
    \bottomrule
  \end{tabularx}
\end{table}


\begin{table}[ht]
  \footnotesize
  \caption{Class~I city and town paradigms.}
  \ifparadigmappendix\else\label{tab:city-town-paradigms}\fi
  \begin{tabularx}{\textwidth}{l l L L L L}
    \toprule
    {\sc noun} & {\sc case} & {\sc sg} & {\sc pl}
      & {\sc sg.def} & {\sc pl.def} \\
    \midrule
    \ird{nuž} & \textsc{dir} & \ird{nuž} & \ird{núže} & \ird{nužo} & \ird{núžete} \\
     & \textsc{trs} & \ird{nuža} & \ird{núžov} & \ird{nužato} & \ird{núžove} \\
     & \textsc{ind} & \ird{nužam} & \ird{núžac} & \ird{nužamo} & \ird{núžace} \\
    \addlinespace
    \ird{vruh} & \textsc{dir} & \ird{vruh} & \ird{vrúhe} & \ird{vruho} & \ird{vrúhete} \\
     & \textsc{trs} & \ird{vruha} & \ird{vrúhov} & \ird{vruhato} & \ird{vrúhove} \\
     & \textsc{ind} & \ird{vruham} & \ird{vrúhac} & \ird{vruhamo} & \ird{vrúhace} \\
    \bottomrule
  \end{tabularx}
\end{table}


\begin{table}[ht]
  \footnotesize
  \caption{Complete paradigms of Classes IV--VI.}
  \ifparadigmappendix\else\label{tab:nm-paradigms-iv-vi}\fi
  \begin{tabularx}{\textwidth}{l l L L L L}
    \toprule
    {\sc noun (cl.)} & {\sc case} & {\sc sg} & {\sc pl}
      & {\sc sg.def} & {\sc pl.def} \\
    \midrule
    \ird{red} (IV) & \textsc{dir} & \ird{red} & \ird{redě} & \ird{redo} & \ird{reďé} \\
     & \textsc{trs} & \ird{rezy} & \ird{redóv} & \ird{rezyte} & \ird{redóve} \\
     & \textsc{ind} & \ird{redě} & \ird{redím} & \ird{reďé} & \ird{redíme} \\
    \addlinespace
    \ird{rán} (V) & \textsc{dir} & \ird{rán} & \ird{rána} & \ird{ráno} & \ird{ránato} \\
     & \textsc{trs} & \ird{ránou} & \ird{rán} & \ird{ránouto} & \ird{ráne} \\
     & \textsc{ind} & \ird{ráne} & \ird{ránom} & \ird{ránete} & \ird{ránome} \\
    \addlinespace
    \ird{pět} (VI) & \textsc{dir} & \ird{pět} & \ird{pětove} & \ird{pěto} & \ird{pětové} \\
     & \textsc{trs} & \ird{pětim} & \ird{pětóv} & \ird{pětimo} & \ird{pětóve} \\
     & \textsc{ind} & \ird{pěty} & \ird{pětovim} & \ird{pětyte} & \ird{pětovime} \\
    \bottomrule
  \end{tabularx}
\end{table}


\begin{table}[ht]
  \footnotesize
  \caption{Complete paradigms of Classes VII--IX.}
  \ifparadigmappendix\else\label{tab:nm-paradigms-vii-ix}\fi
  \begin{tabularx}{\textwidth}{l l L L L L}
    \toprule
    {\sc noun (cl.)} & {\sc case} & {\sc sg} & {\sc pl}
      & {\sc sg.def} & {\sc pl.def} \\
    \midrule
    \ird{tězeň} (VII) & \textsc{dir} & \ird{tězeň} & \ird{tězňa} & \ird{tězno} & \ird{tězňate} \\
     & \textsc{trs} & \ird{těznu} & \ird{těznóv} & \ird{těznuto} & \ird{těznóve} \\
     & \textsc{ind} & \ird{tězně} & \ird{těznem} & \ird{tězňé} & \ird{tězneme} \\
    \addlinespace
    \ird{obel} (VII) & \textsc{dir} & \ird{obel} & \ird{obľa} & \ird{oblo} & \ird{obľate} \\
     & \textsc{trs} & \ird{oblu} & \ird{oblóv} & \ird{obluto} & \ird{oblóve} \\
     & \textsc{ind} & \ird{oblě} & \ird{oblem} & \ird{obľé} & \ird{obleme} \\
    \addlinespace
    \ird{pivo} (VIII) & \textsc{dir} & \ird{pivo} & \ird{piva} & \ird{pivó} & \ird{pivá} \\
     & \textsc{trs} & \ird{piva} & \ird{pivóv} & \ird{pivato} & \ird{pivóve} \\
     & \textsc{ind} & \ird{pivum} & \ird{pivem} & \ird{pivumo} & \ird{piveme} \\
    \addlinespace
    \ird{mamka} (IX) & \textsc{dir} & \ird{mamka} & \ird{mamce} & \ird{mamká} & \ird{mamcé} \\
     & \textsc{trs} & \ird{mamky} & \ird{mamkóv} & \ird{mamkyte} & \ird{mamkóve} \\
     & \textsc{ind} & \ird{mamce} & \ird{mamcám} & \ird{mamcé} & \ird{mamcáme} \\
    \bottomrule
  \end{tabularx}
\end{table}


\begin{table}[ht]
  \footnotesize
  \caption{Complete paradigms of lexically irregular nouns. Direct-plural variants of \ird{študent} are listed on separate lines.}
  \ifparadigmappendix\else\label{tab:nm-irregular-paradigms}\fi
  \begin{tabularx}{\textwidth}{l l L L L L}
    \toprule
    {\sc noun (cl.)} & {\sc case} & {\sc sg} & {\sc pl}
      & {\sc sg.def} & {\sc pl.def} \\
    \midrule
    \ird{pres} (I) & \textsc{dir} & \ird{pres} & \ird{préše} & \ird{preso} & \ird{préšete} \\
     & \textsc{trs} & \ird{presa} & \ird{préšov} & \ird{presato} & \ird{préšove} \\
     & \textsc{ind} & \ird{presam} & \ird{préšac} & \ird{presamo} & \ird{préšace} \\
    \addlinespace
    \ird{naz} (I) & \textsc{dir} & \ird{naz} & \ird{náže} & \ird{nazo} & \ird{nážete} \\
     & \textsc{trs} & \ird{naza} & \ird{nážov} & \ird{nazato} & \ird{nážove} \\
     & \textsc{ind} & \ird{nazam} & \ird{nážac} & \ird{nazamo} & \ird{nážace} \\
    \addlinespace
    \ird{tóm} (I) & \textsc{dir} & \ird{tóm} & \ird{tome} & \ird{tómo} & \ird{tomete} \\
     & \textsc{trs} & \ird{toma} & \ird{tomóv} & \ird{tomato} & \ird{tomóve} \\
     & \textsc{ind} & \ird{tomem} & \ird{tomac} & \ird{tomemo} & \ird{tomace} \\
    \addlinespace
    \ird{študent} (mixed) & \textsc{dir} & \ird{študent} & \shortstack[l]{\ird{študentí}\\\ird{študencí}} & \ird{študento} & \shortstack[l]{\ird{študentíte}\\\ird{študencíte}} \\
     & \textsc{trs} & \ird{študenta} & \ird{študentóv} & \ird{študentato} & \ird{študentóve} \\
     & \textsc{ind} & \ird{študencí} & \ird{študencím} & \ird{študencíte} & \ird{študencíme} \\
    \bottomrule
  \end{tabularx}
\end{table}


\FloatBarrier
\section{Verbal paradigms}

The principal treatments are given in Tables~\ref{tab:verbal-quantity-classes}, \ref{tab:aspect}, \ref{tab:reduced-asp}, and~\ref{tab:inner-evidential-paradigm}.

\begin{table}[htbp]
\footnotesize
\caption{Principal parts of the three verbal quantitative classes.}
\ifparadigmappendix\else\label{tab:verbal-quantity-classes}\fi
\begin{tabularx}{\textwidth}{L L L L L L}
\toprule
{\sc class/stem} & {\sc continuous} & {\sc progressive} & {\sc contemplative} & {\sc perfective} & {\sc retrospective} \\
\midrule
A \ird{zemac-} & \ird{zemaca} & \ird{zemácime} & \ird{zemacáž} & \ird{zemácek} & \ird{zemacáv} \\
B \ird{pjašt-} & \ird{pjašta} & \ird{pjaštíme} & \ird{pjaštáž} & \ird{pjaštník} & \ird{pjaštnáv} \\
C \ird{knář-} & \ird{knářa} & \ird{knářime} & \ird{knářaž} & \ird{knářek} & \ird{knářav} \\
\bottomrule
\end{tabularx}
\end{table}


\begin{table}[ht]
	\footnotesize
	\caption{Aspect-alignment exponents by domain and quantitative class.}\ifparadigmappendix\else\label{tab:aspect}\fi
	\begin{tabularx}{0.98\textwidth}{L L L L}
		\toprule
		{\sc domain} & {\sc aspect} & {\sc exponent} & {\sc quantity and orientation} \\
		\midrule
		Perfective & Perfective & \ird{-ek}/\ird{-ék} & bare forms require an intransitive reading; ordinary transitive \ird{-nik}/\ird{-ník}; class-selected short or long bound \ird{-ik-}/\ird{-ík-} \\
		& Retrospective & \ird{-áv}/\ird{-av} & ordinary A/B long \ird{-áv}, or \ird{-náv} with orientation; C and listed strong-valency stems short \ird{-av}/\ird{-nav} \\
		\midrule
		Imperfective & Continuous & \ird{-a} & local G0 in ordinary A/B; previously selected lexical or valency strength retained \\
		& Progressive & \ird{-ime}/\ird{-íme} & A root G1 with short ending; B suffix-strong; C and strong-valency stem with short ending \\
		& Contemplative & \ird{-áž}/\ird{-až} & A/B ordinary G0 with long ending; C and strong-valency stem with short ending \\
		\bottomrule
	\end{tabularx}
\end{table}


\begin{table}
  \footnotesize
  \caption{Reduced Asp.Align exponents in non-indicative clauses.}
  \ifparadigmappendix\else\label{tab:reduced-asp}\fi
  \begin{tabularx}{0.9\textwidth}{l l L}
    \toprule
    \textsc{domain} & \textsc{exponent} & \textsc{semantic range} \\
    \midrule
    Imperfective (IPFV) & $\varnothing$ & unbounded, ongoing, habitual, generic;
                                        processual commands and optatives \\
    Perfective (PFV)    & \ird{-ov-}/\ird{-óv-} & bounded, single whole event, completive;
                                        ingressively directive \\
    \bottomrule
  \end{tabularx}
\end{table}


\begin{table}[htbp]
\footnotesize
\caption{Indirect evidential forms across the five indicative aspects.}
\ifparadigmappendix\else\label{tab:inner-evidential-paradigm}\fi
\begin{tabularx}{\textwidth}{L L L L}
\toprule
{\sc aspect} & {\sc finite host} & {\sc indirect evidential} & {\sc linked evidential} \\
\midrule
Perfective & \ird{pjaštník} & \ird{pjaštníkal} & \ird{pjaštníkalě} \\
Retrospective & \ird{pjaštnáv} & \ird{pjaštnával} & \ird{pjaštnávalě} \\
Continuous & \ird{pjašta} & \ird{pjaštal} & \ird{pjaštalě} \\
Progressive & \ird{pjaštíme} & \ird{pjaštímel} & \ird{pjaštímelě} \\
Contemplative & \ird{pjaštáž} & \ird{pjaštážal} & \ird{pjaštážalě} \\
\bottomrule
\end{tabularx}
\end{table}


\FloatBarrier
\paradigmappendixfalse
